\section*{الفصل \ref{ch:b2:acyl-substitution} --- مشتقات الأحماض الكربوكسيلية}

\begin{solution}{exo:b2:acyl-substitution:1}
\babelsublr{\foreignlanguage{arabic}{الإيثانأميد}~$<$~\foreignlanguage{arabic}{إيثانوات الإيثيل}~$<$~\foreignlanguage{arabic}{أنهيدريد الإيثانويك}~$<$~\foreignlanguage{arabic}{كلوريد الإيثانويل}}.
\end{solution}

\begin{solution}{exo:b2:acyl-substitution:2}
حمض الإيثانويك (+ \ce{HCl})؛ إيثانوات الإيثيل؛ الإيثانأميد (+ كلوريد الأمونيوم، إذ يحتجز المكافئ الثاني
من الأمونيا \ce{HCl})؛ أنهيدريد الإيثانويك (+ \ce{NaCl})؛
\textit{N}،\textit{N}-ثنائي ميثيل إيثانأميد (+ كلوريد ثلاثي إيثيل الأمونيوم).
\end{solution}

\begin{solution}{exo:b2:acyl-substitution:3}
\ce{C6H5COOCH3 + NaOH -> C6H5COONa + CH3OH}. $13.6/136.15 = \qty{0.0999}{mol}$، والكمية نفسها من
\ce{NaOH}: \qty{4.00}{g}.
\end{solution}

\begin{solution}{exo:b2:acyl-substitution:4}
$\gamma$-بيوتيرولاكتون (أوكسولان-2-ون): حلقة خماسية من أربع ذرات كربون وذرة أكسجين واحدة.
\end{solution}

\begin{solution}{exo:b2:acyl-substitution:5}
تُضاف الأمونيا إلى كربون الكربونيل؛ وينتقل بروتون من النيتروجين إلى الأكسجين (أو إلى
المجموعة المغادرة)؛ ثم يغادر الإيثوكسيد (على هيئة إيثانول، بعد انتقال البروتون) وتتكوّن \ce{C=O} من جديد:
الإيثانأميد والإيثانول.
\end{solution}

\begin{solution}{exo:b2:acyl-substitution:6}
تؤسَّل مجموعة \ce{OH} الفينولية (بحفز حمضي): حمض 2-أسيتوكسي بنزويك؛ والناتج الثانوي حمض الإيثانويك.
\end{solution}

\begin{solution}{exo:b2:acyl-substitution:7}
2-ميثيل بروبان-2-ول. تعطي الإضافة الأولى البروبانون، وهو أكثر محبة للإلكترونات من الإستر، فيأخذ
مجموعة الميثيل الثانية فورًا.
\end{solution}

\begin{solution}{exo:b2:acyl-substitution:8}
البيوتان-1،4-ديول: يُختزل الإستر إلى كحولين أوليين، مع انفتاح الحلقة.
\end{solution}

\begin{solution}{exo:b2:acyl-substitution:9}
مكافئ واحد: $x^2/(1 - x)^2 = 4$، $x = 2/3$. ثلاثة مكافئات: $x^2/[(1 - x)(3 - x)] = 4$، $3x^2 - 16x
+ 12 = 0$، $x = 0.90$. وإزالة الماء حال تكوّنه (دين--ستارك) تدفعه أبعد.
\end{solution}

\begin{solution}{exo:b2:acyl-substitution:10}
يعطي \ce{SOCl2} كلوريد البنزويل؛ ومع مكافئين من ثنائي إيثيل أمين (أو مكافئ واحد مع ثلاثي إيثيل أمين)
يعطي الأميد. أما الحمض والأمين المخلوطان مباشرة فيخضعان لتفاعل حمض--قاعدة: بنزوات ثنائي إيثيل
الأمونيوم، الذي لا يعطي الأميد إلا بتسخين شديد.
\end{solution}

\begin{solution}{exo:b2:acyl-substitution:11}
ثلاثة KOH لكل تريولين: $3 \times 56.11/885.4 = \qty{0.190}{g}$ لكل غرام، أي قرينة \omterm{def:b2:acyl-substitution:saponification}{تصبّن} قدرها
190.
\end{solution}

\begin{solution}{exo:b2:acyl-substitution:12}
\ce{NaBH4}: الكيتون وحده، إلى كحول ثانوي. \ce{LiAlH4}: الكيتون إلى الكحول الثانوي،
والإستر إلى كحول أولي (ويحرر الجزء الكحولي)، والأميد إلى أمين.
\end{solution}

\begin{solution}{pb:b2:acyl-substitution:1}
\textbf{1.} حمض الأولييك \ce{C18H34O2}؛ والغليسرول \ce{C3H8O3}.
\textbf{2.} \ce{C57H104O6}: $57 \times 12.011 + 104 \times 1.008 + 6 \times 15.999 = \qty{885.45}{g/mol}$.
\textbf{3.} ثلاث.
\textbf{4.} الرابطة المزدوجة cis تُحدث انثناءً في السلاسل، فتتراص تراصًا رديئًا: قوى بين جزيئية ضعيفة، 
ودرجة انصهار منخفضة.
\textbf{5.} الغليسرول وثلاثة جزيئات من أوليات الصوديوم، أي صابون.
\textbf{6.} خليط من إسترات الميثيل للأحماض الدهنية.
\textbf{7.} \ce{C57H104O6 + 3CH3OH -> C3H8O3 + 3C19H36O2}.
\textbf{8.} أيون الميثوكسيد، المتكوّن من الميثانول والقاعدة، محب للنواة أقوى بكثير من
الميثانول؛ فيهاجم كربونيلات الإستر.
\textbf{9.} كربون الكربونيل لإحدى مجموعات الإستر حاملًا \ce{O-}، و\ce{OCH3} الآتية من الميثوكسيد، 
وأكسجين الغليسريل: رباعي الأوجه.
\textbf{10.} \omterm{def:b2:acyl-substitution:transesterification}{تبادل الأسترة} توازن ثابته قرب 1؛ والميثانول الزائد يزيحه
نحو إسترات الميثيل.
\textbf{11.} ينفصل الغليسرول طبقةً سفلية: وإزالة ناتج تدفع التوازن.
\textbf{12.} لا: فالميثوكسيد يتجدد في كل تبادل.
\textbf{13.} يعادل القاعدة، معطيًا أوليات الصوديوم (صابونًا): فيُستهلك الحفاز.
\textbf{14.} يصبّنها: فالهيدروكسيد (الناتج من الماء والميثوكسيد) يقطع الإسترات إلى صابون 
وميثانول، على نحو غير عكوس.
\textbf{15.} يستهلك القاعدة، ويستحلب الخليط ويمنع طبقة الغليسرول من الانفصال.
\textbf{16.} \omterm{def:b2:acyl-substitution:esterification}{بأسترة} محفَّزة بالحمض مع الميثانول أولًا، تحوّل الأحماض الحرة إلى
إسترات ميثيل.
\textbf{17.} يحلمه الماء الإسترات ويحوّل الحفاز إلى هيدروكسيد، يقوم \omterm{def:b2:acyl-substitution:saponification}{بالتصبين}.
\textbf{18.} $10^6/885.45 = \qty{1129}{mol}$.
\textbf{19.} $3 \times 1129 \times 32.04 = \qty{108.6}{kg}$.
\textbf{20.} $3 \times 1129 \times 296.50 = \qty{1004.6}{kg}$.
\textbf{21.} الداخل $1000 + 108.6 = \qty{1108.6}{kg}$؛ والخارج $104.0 + 1004.6 = \qty{1108.6}{kg}$.
\textbf{22.} \qty{20}{kg} من حمض الأولييك، أي $20\,000/282.47 = \qty{70.8}{mol}$، والكمية نفسها من \ce{NaOH}:
\qty{2.83}{kg}.
\textbf{23.} $1129 \times 92.09 = \boldsymbol{\approx \qty{104}{kg}}$ من الغليسرول لكل طن.
\end{solution}
